\documentclass[11pt]{article}
% A two-page summary of the Ontologizer for readers new to the project.
% Kept in plain language on purpose: every term is defined where it first
% appears, and project vocabulary from findings.tex is avoided. Numbers
% come from findings.tex; change them there first.
\usepackage[biblatex]{flockpaper}
\usepackage[backend=biber,style=numeric,sorting=none]{biblatex}
\addbibresource{refs.bib}

\title{Unsupervised Steering Vectors from Learned Classifiers:\\
A Two-Page Summary}
\author{%
  \begin{tabular}[t]{@{}l@{\hspace{2em}}l@{}}
    Keira Wiechecki & Lamb Wakefield \\[2pt]
    \mutedlabel{Wakaru \textperiodcentered{} kaw504@nyu.edu} &
    \mutedlabel{Flock Capital \textperiodcentered{} lamb@flock.holdings} \\[9pt]
    Jade Zaslavsky & vmfunc \\[2pt]
    \mutedlabel{Wakaru \textperiodcentered{} jadeofnameless@gmail} &
    \mutedlabel{Flock Capital \textperiodcentered{} celeste@linux.com}
  \end{tabular}}
\date{\today}

\begin{document}
\maketitle

\section*{What the method does}

The method, which we call an \emph{Ontologizer}, learns a set of small
classifiers that together describe a neural network's internal
representation. Each classifier sorts every input into one of $k$
categories (we use $k=32$), and each category stores one vector. An input
is described by which category each classifier chose, and it is
reconstructed as the sum of the chosen vectors. Nothing is labelled: the
classifiers and their vectors are learned only by asking the sum to
reproduce the representation as closely as possible.

Because each classifier's categories are mutually exclusive, changing one
classifier's answer for an input --- from category $a$ to category $b$ ---
swaps one stored vector for another. That swap is a steering intervention:
it moves the representation from ``$a$'' to ``$b$'' along a single,
named axis, and the model supplies the direction.

\section*{The problem it addresses}

To steer a model is to change its internal state so that it behaves
differently in a chosen way. The most reliable current method is a
\emph{steering vector}: take the average representation of examples with a
property, subtract the average of examples without it, and add the
difference to new inputs~\cite{turner2023,rimsky2024}. This works well,
but someone must decide which property to look for and supply labelled
examples of both sides, so it can only find what its user already
suspects is there.

Sparse autoencoders (SAEs) find features without labels: each input is
described by a few active features out of thousands~\cite{bricken2023,
gao2024}. But their features switch on and off independently, and adding
one to an input tends to change others and to damage the input in
unpredictable ways; on standard steering benchmarks, simple supervised
baselines beat them~\cite{wu2025axbenchsteeringllmssimple}. The goal here
is steering vectors that need no labels but keep the useful shape of the
supervised kind: each one a choice between alternatives on a shared axis.

\section*{What it assumes, and why}

\begin{enumerate}
\item \textbf{Concepts are directions that add.} A representation is,
approximately, a sum of vectors for the things it encodes. This is the
linear representation hypothesis~\cite{park2023linear}, and it is what
makes supervised steering vectors work at all.
\item \textbf{Much of the variation is categorical.} A sentence has one
dominant topic, one language, one register; it is more natural to describe
it by choosing among alternatives than by switching many independent
features on and off. We check this rather than assume it: the clearest
classifier the method learns on sentence embeddings sorts by topic and
genre (crime news, legislation, encyclopedic description, and so on).
\item \textbf{The categories are a property of the data, not of the
training run.} If the method recovers real structure, a second run from a
different random start should find the same classifiers. This is the
assumption that has \emph{not} held, and it is the main open problem
(below).
\end{enumerate}

\section*{How it works, intuitively}

Think of describing a sentence by answering a fixed list of multiple-choice
questions, where the questions and answers are learned. The first round of
questions captures the broad outline; the model then measures what that
outline gets wrong and asks a second round about the remaining error, and
so on for a few rounds. The answers can be soft (a mixture of categories)
or hard (exactly one category per question). In the hard version an input
is stored as a short list of discrete choices --- 1900 bits for a sentence
embedding --- with no continuous numbers at all.

We apply it to two representations: SONAR sentence
embeddings~\cite{sonar}, which can be decoded back into text, so a steered
embedding can be read; and the internal activations of GPT-2
small~\cite{radford2019}, where the reconstruction can be put back into
the model to see how much of its behaviour survives.

\section*{How it compares with existing methods}

\textbf{Against supervised steering --- the comparison that matters.} We
have one direct test. Steering sentence embeddings toward a target
language, forcing the classifier the method learned moved the decoded text
toward that language $0.43$ times as far as the supervised
difference-of-means vector did, and 5 times further than a linear-probe
direction. The method never saw a label in training; labels were used only
to pick which category to force. All methods, the supervised one included,
damaged unrelated content about as much as a random step of the same size.
So on this test the method recovers a useful but weaker steering
direction than supervision provides. It has \emph{not} yet been shown to
beat supervised steering.

\textbf{Against SAEs on steering.} Measured directly in embedding space,
at each model's own natural step size, a feature of a standard SAE (32
active features per input) takes effect on 46--50\% of inputs, and other
SAE settings range from 8\% to 82\%; the method's multi-round version takes
effect on only 8--19\%, while disturbing other classifiers as much as a
random step would.
A simplified one-round version of the method does best of all on side
effects: it takes effect on 41\% of inputs while disturbing only 1.4\% of
the other classifiers, against 18\% for a random step.

\textbf{Against SAEs on reconstruction (a secondary check).} Describing a
sentence embedding with 1900 bits of discrete choices, the method leaves
$15\%$ of the variance unexplained, against $26\%$ for an SAE using a
similar number of bits plus 160 continuous coefficients. On GPT-2, an
800-bit code preserves $97.6\%$ of the model's performance, against
$96.2\%$ for an SAE with similar error.

\section*{What does not work yet}

\begin{itemize}
\item \textbf{The classifiers are not reproducible.} Two runs from
different random starts reconstruct equally well, but agree on almost none
of their classifiers beyond the first round; only the topic classifier
reappears reliably. A classifier that is different on every run cannot be
named, so this limits everything else.
\item \textbf{Steering through the input is not targeted.} Pushing an input
toward one category changes other classifiers' answers about as much as a
random push of the same size, and the multi-round structure doubles that.
\item \textbf{The stored vectors cannot be read on their own.} Decoding a
category's vector into text does not describe what the category selects;
examples of inputs are needed.
\end{itemize}

\section*{Next}

The test the method needs is a head-to-head steering benchmark:
difference-of-means vectors, probe directions, SAE features and the
method's own swaps, scored on the same targets with the same success and
side-effect measures at matched step sizes. That comparison is being set
up now. In parallel, the reproducibility problem has to be solved, or the
reproducible classifiers identified, before any classifier is given a
name.

\printbibliography

\end{document}
